%% ma: L12-B14-0128
%% lop: 12
%% chuong: 5
%% bai: 14
%% dang: 12.14.1
%% loai: TN4
%% muc: TH
%% dap_an: "B"
%% nguon: "Ảnh giáo viên gửi 25/11/2026 (tài liệu ôn tập Phương trình mặt phẳng), Phần I Câu 6"
%% kien_thuc: "Bài 14 (hai vectơ pháp tuyến cùng phương)"
%% trang_thai: chinh-thuc
\begin{ex}
Trong không gian $Oxyz$, cho mặt phẳng $(P)$ có phương trình $2x+3y-z+4=0$. Biết $\vec n=(1;b;c)$ là một vectơ pháp tuyến của $(P)$. Khi đó, tổng $T=b+c$ bằng
\choice{$2$}{\True $1$}{$4$}{$0$}
\loigiai{$(P)$ có vectơ pháp tuyến $(2;3;-1)$. Vì $\vec n=(1;b;c)$ cùng phương với $(2;3;-1)$ nên $\vec n=\dfrac12(2;3;-1)=\left(1;\dfrac32;-\dfrac12\right)$, suy ra $b=\dfrac32$, $c=-\dfrac12$ và $T=1$.}
\end{ex}
